\documentclass[11pt]{article}

\usepackage[T1]{fontenc}
\usepackage[utf8]{inputenc}
\usepackage{amsmath,amssymb,amsthm}
\usepackage{booktabs}
\usepackage{geometry}
\usepackage{hyperref}
\usepackage{microtype}
\usepackage{xcolor}

\geometry{margin=1in}
\setlength{\emergencystretch}{2em}
\hypersetup{
  colorlinks=true,
  linkcolor=blue!55!black,
  citecolor=blue!55!black,
  urlcolor=blue!55!black
}

\newtheorem{theorem}{Theorem}[section]
\newtheorem{lemma}[theorem]{Lemma}
\newtheorem{proposition}[theorem]{Proposition}
\newtheorem{corollary}[theorem]{Corollary}
\theoremstyle{definition}
\newtheorem{definition}[theorem]{Definition}
\theoremstyle{remark}
\newtheorem{remark}[theorem]{Remark}

\newcommand{\F}{\mathbb F}
\newcommand{\A}{\mathcal A}
\newcommand{\FiveCDC}{\textup{FiveCDC}}
\newcommand{\Span}{\operatorname{span}}
\numberwithin{equation}{section}

\title{Six-Point One-Hole Compression for\\
Five-Cycle Double Covers}
\author{Atharva Vaidya\\
\small AI-assisted research draft for independent human verification}
\date{July 29, 2026}

\begin{document}
\maketitle

\begin{center}
\fcolorbox{red!65!black}{red!4}{
\parbox{0.91\textwidth}{
\textbf{Resolution status and trust boundary.}
The standard five-cycle double cover conjecture remains open.  This
manuscript proves an exact affine characterization after a nowhere-zero
\(\F_2^3\)-flow, an omitted coordinate pair, and a merge pair have been
fixed.  Existentially quantifying those choices gives exactly the original
conjecture, not a proof of it.  The orientable conjecture is not considered.

OpenAI Codex agents, under Atharva Vaidya's direction, discovered the
normal form, derived the proof and local tables, wrote both computational
implementations, performed the enumerations, and drafted this manuscript.
The second implementation is independent of the first implementation, not
independent of AI.  No human peer review is claimed.  A human author must
check the proof, rerun or reimplement the finite audits, review prior art,
and accept scholarly responsibility before submission.
}}
\end{center}

\begin{abstract}
Let \(W=\F_2^3\), let \(G\) be a finite loopless cubic multigraph, and
fix a nowhere-zero \(W\)-flow \(f\).  The standard eight-coordinate
triangle construction labels every edge by a two-subset of \(W\).  We
show that this construction compresses to a five-cycle double cover if
its labels use only six points and omit one pair among them: identify the
two points of that unused pair.

For fixed disjoint pairs \(R,M\in\binom W2\), where \(R\) is omitted
from the support and \(M\) is the pair to be identified, existence of
such a labelling is exactly an affine system over \(\F_2\), with three
potential bits per graph vertex.  We give a complete human proof and
classify every local list.  Up to affine equivalence, \(R\) and \(M\)
are parallel or skew; the seven local list sizes are respectively
\[
              (4,2,2,2,2,2,2)
       \quad\hbox{and}\quad
              (4,4,2,2,2,1,1).
\]
Failure for a fixed triple \((f,R,M)\) has the elementary dual
certificate \(y^{\mathsf T}A=0\), \(y^{\mathsf T}b=1\).

Two independently written exact checkers verify the local classification
and a strict 34-vertex fixed-flow example: all 56 five-point restrictions
fail, while eight of the 420 six-point one-hole systems succeed.  A
complete 900-orbit flow census on one rigid 12-vertex graph separates the
five-point, six-point one-hole, and unrestricted five-colourable
coordinate-potential tests.  These are structural and computational
results about fixed flows.  They neither prove nor refute \FiveCDC.
\end{abstract}

\section{Definitions and exact scope}

A \emph{cycle} in this manuscript is an Eulerian edge-subset: every
vertex has even degree in it.  It need not be connected or 2-regular.
A \emph{five-cycle double cover} is an indexed family
\((C_1,\ldots,C_5)\) of such subsets in which every edge occurs exactly
twice.  Empty members are allowed.  Thus ``five'' means ``at most five,''
as in the standard formulation used by Oum~\cite{Oum2026}.  Parallel
edges are distinct edge-set elements.  We exclude loops here and make no
claim about the orientable variant.

The standard conjecture, attributed independently to Celmins and
Preissmann, asserts that every finite bridgeless graph has a five-cycle
double cover~\cite{Celmins1984,Preissmann1981,HusekSamal2026,Oum2026}.
It remains open.  Hu\v{s}ek and \v{S}\'amal recently gave an equivalent
flow-selection condition for cubic graphs
\cite[Theorem~3.16 and Conjecture~3.19]{HusekSamal2026}.  Our result is
another exact fixed-flow normal form.  Its graph-level existential closure
is again \FiveCDC, so the remaining quantifier cannot be omitted.

There is important earlier configuration-colouring literature.
Kr\'al', M\'a\v{c}ajov\'a, Pangr\'ac, Raspaud, Sereni, and
\v{S}koviera developed projective, affine, and abelian colourings of
cubic graphs, including Desargues-configuration equivalences and the
\(F_i/A_i\) colouring framework~\cite{KralEtAl2009}.  We claim no novelty
for those configuration or colouring equivalences.  The possible new
contribution here is narrower: after fixing a Fano flow, we impose a
six-point support and one missing duad, derive its affine potential system
and complete local orbit table, and certify a strict separation from the
five-point restriction for one fixed flow.  Hu\v{s}ek--\v{S}\'amal are
also the prior source used here for the eight-coordinate system and its
five-set/component-parity specialization~\cite{HusekSamal2026}.

Put \(W=\F_2^3\), written additively.  A map
\[
                    f:E(G)\longrightarrow W-\{0\}
\]
is a nowhere-zero \(W\)-flow when
\[
                    \sum_{e\ni v}f(e)=0
                    \qquad(v\in V(G)).
\]
At a cubic vertex the three incident values are the nonzero elements of
a unique two-dimensional linear subspace \(H_v\le W\).

\begin{definition}
An \emph{eight-coordinate pair labelling} is a map
\[
                    P:E(G)\longrightarrow \binom W2
\]
such that, for every \(v\in V(G)\) and \(z\in W\), the number of
incident edges \(e\ni v\) with \(z\in P_e\) is even.  It is
\emph{compatible with \(f\)} when the xor of the two members of \(P_e\)
is \(f(e)\) for every edge \(e\).
\end{definition}

Every coordinate \(z\) defines the Eulerian set
\(\{e:z\in P_e\}\), and every edge belongs to exactly two coordinates.
Thus a pair labelling is an eight-cycle double cover with named
coordinates.  We next parameterize all labellings compatible with a
fixed flow.

\section{Coordinate triangles and vertex potentials}

For a two-dimensional subspace \(H\le W\) and \(t\in W\), define
\begin{equation}
                         \Delta_H(t)=(t+H)-\{t\}.
                         \label{eq:triangle}
\end{equation}
This is a three-point affine triangle.  Its three sides have differences
equal to the three nonzero elements of \(H\).

For an incidence \(e\ni v\), put \(p=f(e)\), and let \(s_{v,e}\) be
either one of the other two flow values at \(v\).  The side of
\(\Delta_{H_v}(t_v)\) having difference \(p\) is
\begin{equation}
 P_{v,e}
 =t_v+s_{v,e}+\langle p\rangle
 =\{t_v+s_{v,e},\,t_v+s_{v,e}+p\}.
 \label{eq:side}
\end{equation}
The two possible choices of \(s_{v,e}\) differ by \(p\), so the unordered
pair in \eqref{eq:side} is well defined.

\begin{lemma}[potential parameterization]
\label{lem:potentials}
For a fixed nowhere-zero flow \(f\), compatible pair labellings are in
bijection with vertex potentials \((t_v:v\in V(G))\) satisfying
\begin{equation}
 t_u+t_v+s_{u,e}+s_{v,e}\in\langle f(e)\rangle
 \qquad(e=uv).
 \label{eq:compatibility}
\end{equation}
The labelling is then given by \eqref{eq:side}.
\end{lemma}

\begin{proof}
The two endpoint formulae in \eqref{eq:side} describe the same unordered
pair exactly when their coset representatives differ by an element of
\(\langle f(e)\rangle\), which is \eqref{eq:compatibility}.  At a vertex,
the three pairs are the three sides of \(\Delta_{H_v}(t_v)\), so each of
its three coordinate points occurs twice and every other point occurs
zero times.  The coordinate parity condition follows.

Conversely, at a cubic vertex, the three labels of any pair labelling
form a triangle.  Indeed, if the first two pairs are \(A,B\), parity
forces the third to be \(A\mathbin{\triangle}B\); since it has size two,
\(|A\cap B|=1\).  Write the resulting triangle points as \(x,y,z\).
Its edge differences are the three incident flow values, and
\[
                              t_v=x+y+z
\]
is the unique potential for which its point set is
\(\Delta_{H_v}(t_v)\).  Agreement of each edge label at its endpoints
then gives \eqref{eq:compatibility}.  The two constructions are inverse.
\end{proof}

For any nonzero \(p\in W\), membership in \(\langle p\rangle\) is
equivalent to two scalar affine equations: apply a basis of the
two-dimensional annihilator of \(p\) in \(W^*\).  Hence all edge
conditions in \eqref{eq:compatibility} are affine binary equations.

\section{The six-point one-hole system}

Choose disjoint pairs
\begin{equation}
                         R,M\in\binom W2,\qquad R\cap M=\varnothing.
                         \label{eq:pairs}
\end{equation}
The points of \(R\) are forbidden, so the allowed support is
\(S=W-R\), of size six.  The pair \(M\subset S\) is also forbidden as
an edge label.  For a Fano plane \(H\), define
\begin{equation}
 \A_H(R,M)=
 \{t\in W:\Delta_H(t)\subseteq W-R,
                 M\nsubseteq\Delta_H(t)\}.
 \label{eq:local-list}
\end{equation}
The second condition says precisely that \(M\) is not one of the three
sides of the local triangle.

\begin{theorem}[affine six-point one-hole criterion]
\label{thm:affine}
Fix \(G,f,R,M\) as above.  There is a compatible pair labelling supported
on \(W-R\) and not using the pair \(M\) if and only if the following
system is soluble:
\begin{align}
 t_v&\in\A_{H_v}(R,M)
                         &&(v\in V(G)), \label{eq:local-system}\\
 t_u+t_v+s_{u,e}+s_{v,e}
     &\in\langle f(e)\rangle
                         &&(e=uv). \label{eq:edge-system}
\end{align}
Every set \(\A_H(R,M)\) is a nonempty affine subspace of \(W\), of
dimension zero, one, or two.  Consequently
\eqref{eq:local-system}--\eqref{eq:edge-system} is an ordinary affine
system over \(\F_2\).
\end{theorem}

\begin{proof}
By Lemma~\ref{lem:potentials}, condition \eqref{eq:edge-system} is
equivalent to endpoint compatibility.  The definition
\eqref{eq:local-list} is exactly the support and missing-pair condition
at each vertex.  It remains only to prove that every local list is affine
and nonempty.  This follows from the exhaustive two-case calculation in
Theorem~\ref{thm:local-table} below.
\end{proof}

To turn a local membership condition into scalar equations, write its
nonempty affine subspace as \(t_0+U\).  Then
\[
        t\in t_0+U
        \quad\Longleftrightarrow\quad
        \lambda(t)=\lambda(t_0)
        \quad\hbox{for every }\lambda\in U^\perp.
\]
Together with the two annihilator equations per edge, this gives a
literal matrix \(Ax=b\) with three bits per vertex.

\section{Complete local classification}

Translate the eight coordinate points so that
\begin{equation}
                    R=\{0,d\},\qquad M=\{a,b\},\qquad m=a+b.
                    \label{eq:normal}
\end{equation}
Translation changes neither pair differences nor the planes \(H\).
There are two affine orbits:
\begin{itemize}
\item \emph{parallel}, when \(m=d\); then
      \(R\cup M=H_0:=\langle d,a\rangle\);
\item \emph{skew}, when \(m\ne d\); disjointness implies that
      \(d,m,a\) form a basis of \(W\).
\end{itemize}

\begin{theorem}[local list table]
\label{thm:local-table}
In the parallel case,
\begin{equation}
\A_H(R,M)=
\begin{cases}
W-H,&H=H_0,\\
M,&d\in H,\ H\ne H_0,\\
R,&d\notin H.
\end{cases}
\label{eq:parallel-table}
\end{equation}
Across the seven planes, the list sizes are
\[
                         (4,2,2,2,2,2,2).
\]
In the skew case,
\begin{equation}
\A_H(R,M)=
\begin{cases}
M,&d,m\in H,\\
W-H,&d\in H,\ m\notin H,\\
\{r\in R:a+r\notin H\},&d\notin H,\ m\in H,\\
R,&d,m\notin H.
\end{cases}
\label{eq:skew-table}
\end{equation}
Across the seven planes, the list sizes are
\[
                         (4,4,2,2,2,1,1).
\]
In particular, every local list is a point, affine line, or affine
plane.
\end{theorem}

\par\medskip
\begin{proof}
First ignore the restriction on \(M\).  If \(d\in H\), the forbidden
points \(0,d\) lie in one \(H\)-coset.  The triangle must therefore use
the other coset, and its omitted point \(t\) may initially be any of the
four points of that coset.  If \(d\notin H\), the forbidden points lie in
opposite \(H\)-cosets, and the only possible omitted points are
\(t=0,d\).

Now impose the missing-pair condition.  If \(m\notin H\), the two members
of \(M\) lie in opposite \(H\)-cosets, so they cannot both occur in one
triangle.  If \(m\in H\), they lie in one coset.  When that coset is the
chosen one, one of \(a,b\) must itself be the omitted point.

In the parallel case \(m=d\).  For \(H=H_0\), the allowed potentials are
the other coset \(W-H\).  For the other two planes containing \(d\), the
missing-pair condition leaves \(M\).  For each of the remaining four
planes it leaves \(R\).  This is \eqref{eq:parallel-table}.

In the skew case, exactly one plane contains both \(d,m\); its list is
\(M\).  Two planes contain \(d\) but not \(m\), giving \(W-H\).  Two
contain \(m\) but not \(d\); the initially allowed set \(R\) loses exactly
one point according as \(a+r\in H\).  The final two planes contain
neither \(d\) nor \(m\), and their list is \(R\).  This is
\eqref{eq:skew-table} and gives the stated sizes.
\end{proof}

There are \(28\cdot15=420\) ordered disjoint choices of \((R,M)\):
84 parallel and 336 skew.  The sum of either seven-entry profile is 16.
Thus the table itself predicts
\[
               84\cdot16+336\cdot16=6{,}720
\]
admitted local triangles among the \(420\cdot7=2{,}940\)
pair-pair-plane cases.  The two supplied programs reproduce these counts
directly.

\section{Compression to five coordinates}

\begin{proposition}[one-hole quotient]
\label{prop:quotient}
Every solution of Theorem~\ref{thm:affine} yields a standard five-cycle
double cover of \(G\).
\end{proposition}

\begin{proof}
Write \(M=\{a,b\}\).  Identify \(a\) and \(b\), while leaving the other
four members of \(S=W-R\) distinct.  This gives five coordinate names.
No used pair collapses to a singleton, because \(M\) is the only pair
whose endpoints are identified and it is unused.

Before identification, every old coordinate has even degree at each
graph vertex.  After identification, the degree of the merged coordinate
is, modulo two, the sum of the old \(a\)- and \(b\)-degrees.  It is
therefore even, as are the four unchanged coordinates.  Each edge remains
labelled by two distinct new names, so it is covered exactly twice.
\end{proof}

The 14 permitted pair types make the compression explicit.  Six pairs
avoid both \(a,b\) and map injectively to six target pairs.  The eight
pairs joining one of \(a,b\) to one of the other four points map
two-to-one to the remaining four target pairs.  This is not a fixed
pointwise map from the seven nonzero flow values to ten target pairs:
two preimages of one target pair have differences separated by \(a+b\).

\begin{corollary}[exact graph-level scope]
\label{cor:equivalence}
A finite loopless cubic multigraph has a standard five-cycle double cover
if and only if there exist a nowhere-zero \(W\)-flow \(f\) and disjoint
pairs \(R,M\) for which
\eqref{eq:local-system}--\eqref{eq:edge-system} is soluble.
\end{corollary}

\begin{proof}
The reverse implication is Proposition~\ref{prop:quotient}.  Conversely,
label each edge of a five-cycle double cover by the two cover indices
that contain it.  At a cubic vertex every coordinate occurs zero or two
times, so the three edge labels form a triangle, as in the proof of
Lemma~\ref{lem:potentials}.

Inject the five coordinate names into any five members of a six-set
\(S\subset W\), leaving one point \(q\in S\) unused.  Put \(R=W-S\),
and let \(M\) join \(q\) to any used point.  No original edge uses \(M\).
The xor of the two endpoint names of each edge label is nonzero and sums
to zero at every vertex, hence gives a nowhere-zero \(W\)-flow.  The
original local triangles give potentials solving the affine system.
\end{proof}

\begin{remark}
Corollary~\ref{cor:equivalence} is a reformulation of \FiveCDC, not a
resolution.  The new tractability appears only after \(f,R,M\) are
fixed.  A proof of the conjecture would still have to select a successful
triple universally; a counterexample would have to exclude every triple.
\end{remark}

\section{Dual certificates for a fixed failure}

After scalarization, write the fixed system as
\begin{equation}
                              Ax=b
                              \label{eq:matrix}
\end{equation}
over \(\F_2\).  The variables are the three coordinate bits of every
vertex potential.  Local rows annihilate the direction space of the
appropriate affine list, and two edge rows annihilate each edge value.

\begin{proposition}[left-kernel certificate]
\label{prop:dual}
The fixed system \eqref{eq:matrix} is inconsistent if and only if there
is a binary row multiplier \(y\) such that
\begin{equation}
                    y^{\mathsf T}A=0,\qquad
                    y^{\mathsf T}b=1.
                    \label{eq:dual}
\end{equation}
\end{proposition}

\begin{proof}
If \eqref{eq:dual} holds, xoring the selected equations gives \(0=1\),
so no solution exists.  Conversely, row reduction of an inconsistent
binary affine system produces a zero coefficient row with right side one.
Keeping the same row operations on an identity matrix expresses that row
as a multiplier \(y\) of the original rows.
\end{proof}

This is a short human-checkable certificate for one failed choice
\((f,R,M)\).  It is not a certificate that every choice fails.  The
distinction is essential: a graph-level counterexample would require a
certified universal exclusion over flows as well as pairs.

\section{A strict 34-vertex fixed-flow separation}

The retained simple connected bridgeless cubic graph has graph6 record
\begin{center}
\small
\path{as???SK????A?A?C_@_?B?@_??G??_I?GA_AA????_??@?????B??@_?o??????_C??CGO??B@????a????__??C@O??A?_}
\end{center}
and 51 edges.  Its displayed Fano flow and compatible pair labels are
frozen in the project artifact
\begin{center}
\small
\path{output/jaeger-star-thinning-countermodel-34v/star-good-six-witness.json}
\end{center}
whose SHA-256 is
\begin{center}
\scriptsize
\nolinkurl{e3160e5f6b2985823a3421e9f3b1a4582c9c8be4b9abfd1f0e2312d311eb8358}.
\end{center}
\par\noindent
The literal cover uses the six points \(\{0,3,4,5,6,7\}\) and all 14
pair types except \(\{4,5\}\).  Identifying 4 and 5 gives its displayed
\FiveCDC.

\par\medskip
\begin{proposition}[verified strict separation]
\label{prop:strict}
For this fixed flow:
\begin{enumerate}
\item none of the 56 five-point restrictions is soluble;
\item exactly eight of the 420 six-point one-hole systems are soluble;
\item all eight successful choices are skew and have affine solution
      dimension one.
\end{enumerate}
\end{proposition}

This finite proposition is obtained by exact enumeration in two
independently structured programs.  The primary Python checker uses full
row reduction; the JavaScript checker uses an incremental highest-pivot
xor basis.  Both reconstruct the graph, validate simplicity,
connectedness, cubicity, bridgelessness, the flow, and the literal cover.
They then test all 56 and 420 choices.

Two five-row dual certificates illustrate that individual failures can
also be checked by hand.  For the failed five-set
\(\{0,1,2,3,4\}\), vertices 0 and 1 both have forced potential 7.
At each vertex, take the local rows for functionals 1 and 2; their xor
is functional 3 and their two right sides sum to zero.  Edge 0 joins
the two vertices, has flow value 3, and its functional-3 compatibility
row has right side one.  Xoring the five rows cancels every variable
and yields \(0=1\).

For the failed six-hole choice \(R=\{0,1\}\), \(M=\{2,3\}\), vertices
4 and 8 both have local list \(\{2,3\}\).  At each vertex, xor the
functional-4 row of right side zero with the functional-6 row of right
side one; the coefficient is functional 2.  Edge 5 joins the vertices,
has flow value 1, and its functional-2 compatibility row has right side
one.  Again all coefficients cancel while the five right sides xor to
one.  The audit emits and rechecks the exact row indices and metadata.

The older five-point restriction has local profile
\[
                         (4,1,1,1,1,1,1).
\]
Only one Fano line carries freedom; eliminating forced variables gives
the component-parity system of Hu\v{s}ek and \v{S}\'amal
\cite{HusekSamal2026}.  The profiles in
Theorem~\ref{thm:local-table} leave freedom on several lines, allowing
their global equations to couple.  Proposition~\ref{prop:strict} proves
that this is a genuine enlargement at fixed flow.

\section{A complete all-flow calibration on 12 vertices}

For the cubic graph with graph6 record
\[
                             \texttt{K??FEaKR@oE\_},
\]
a retained complete census contains 900
\(\operatorname{GL}(3,2)\)-orbits of nowhere-zero flows, representing
150,192 labelled flows.  Global translation reduces the 56 five-point
supports to seven representatives and the 420 pairs \((R,M)\) to 63
representatives.

For comparison with the two restricted tests, call a compatible
potential \emph{generally successful} when its coordinate co-occurrence
graph on \(W\) is properly colourable with at most five colours.  Merging
coordinates by such a colouring gives a \FiveCDC.  Conversely, every
\FiveCDC{} embeds in this form; this is the unrestricted compatible-potential
test for a fixed flow.

\begin{proposition}[verified all-flow census]
\label{prop:census}
Two exact implementations independently obtain the following partition.
\begin{center}
\begin{tabular}{@{}lrr@{}}
\toprule
fixed-flow outcome & flow orbits & labelled flows\\
\midrule
five-point soluble & 566 & 94,080\\
five-point insoluble, six-hole soluble & 252 & 42,336\\
six-hole insoluble, generally successful & 18 & 3,024\\
not generally successful & 64 & 10,752\\
\bottomrule
\end{tabular}
\end{center}
\end{proposition}

Thus six-point one-hole compression rescues 252 of the 334 flow orbits
that fail every five-point restriction.  It is strictly broader than the
five-point criterion, but not the most general compatible-potential
compression: 18 further orbits need a co-occurrence pattern using more
than six coordinate points.  The 64 remaining flows are obstructions only
for those fixed flows, not for the graph.

One rigid flow among the last 64 has edge-value word
\[
\begin{split}
 &(1,2,3,4,7,3,5,3,6,5,7,2,5,1,4,6,2,4).
\end{split}
\]
Its gauged compatible potential is unique and its co-occurrence graph is
\(K_6\), using all 15 pairs on six points.  Hence translation cannot
create a missing pair.  The same graph has other flows that pass the
stronger five-point test, demonstrating again why the fixed-flow and
graph-level quantifiers must be separated.

\section{Reproducibility}

The exact artifact package is
\begin{center}
\path{scratch/six-point-one-hole-affine-20260729/}
\end{center}
From the project root, the four principal commands are:
\begin{verbatim}
python3 scratch/six-point-one-hole-affine-20260729/\
audit_six_point_one_hole_affine.py
node scratch/six-point-one-hole-affine-20260729/\
verify_six_point_one_hole_affine.mjs
python3 scratch/six-point-one-hole-affine-20260729/\
audit_all_flows_12v.py
node scratch/six-point-one-hole-affine-20260729/\
verify_all_flows_12v.mjs
\end{verbatim}
The line continuations above are typographical; the repository README
gives literal one-line shell commands.  All four reports end with status
\texttt{PASS}.  The package's SHA-256 ledger has digest
\begin{center}
\scriptsize
\nolinkurl{8e89160b7eac0f7aaf5c4fe1b369bddbf89ae7f013211b4022257dfc09258349}.
\end{center}
The first two programs derive the local table, inspect all 2,940 local
cases, and replay Proposition~\ref{prop:strict}; the latter two independently
replay Proposition~\ref{prop:census}.  The reference-hash ledger also
freezes every external input used by these checks.

The human proof of Theorems~\ref{thm:affine} and
\ref{thm:local-table}, Proposition~\ref{prop:quotient}, and
Corollary~\ref{cor:equivalence} is entirely contained in this manuscript.
The strict witness and census propositions depend on exact finite
enumeration.  No SAT solver or uncheckable UNSAT assertion is used in
those two propositions.

\section{What is new, and what remains open}

The conservative possible-novelty claim is the combination of:
\begin{enumerate}
\item the six-point one-hole affine normal form for fixed
      \((f,R,M)\);
\item the complete parallel/skew local-list classification;
\item the left-kernel certificate interface for a failed fixed choice;
\item the strict 34-vertex fixed-flow separation from every five-point
      restriction; and
\item the exact 12-vertex all-flow calibration.
\end{enumerate}
The project has not established priority.  The relevant literature on
Fano flows, generalized colourings, and the new component-parity
criterion includes Jaeger~\cite{Jaeger1979}, the projective/affine/abelian
colouring framework of Kr\'al' et al.~\cite{KralEtAl2009},
Oum~\cite{Oum2026}, and Hu\v{s}ek--\v{S}\'amal
\cite{HusekSamal2026}.  In particular, no novelty is claimed for the
Desargues-configuration equivalence, \(F_i/A_i\) colourings, the
eight-coordinate construction, or the five-point component-parity
criterion.  A specialist literature review is required before asserting
novelty in a submitted version.

The universal remaining statement is: for every graph in a sound
reduced cubic class, select some nowhere-zero \(W\)-flow \(f\) and some
disjoint \(R,M\) for which the affine system is soluble.  By
Corollary~\ref{cor:equivalence}, this is exactly \FiveCDC.  Complete
finite success of selected flow families is evidence, not an induction
theorem.  Conversely, failure of one or many fixed systems is not a graph
counterexample.  No such counterexample, and no universal selection
proof, is supplied here.

\section*{AI-use disclosure}

OpenAI Codex agents using GPT-5-series models, under Atharva Vaidya's
direction, proposed the six-point one-hole compression, derived the
affine equations and local orbit tables, found and interpreted the finite
examples, wrote the Python and JavaScript implementations, executed the
computations, and drafted essentially all prose and \LaTeX{} in this
manuscript.  Separate agents performed adversarial checks, but they were
also AI systems.  Atharva Vaidya directed the project and is responsible
for deciding whether and how to disseminate the work.

The paper deliberately distinguishes human-checkable algebra from
finite computation.  Reproducibility and independently structured code
do not constitute independent human verification.  No cited researcher
is represented as endorsing these results.  Before public submission, a
human author should verify the displayed proofs line by line, inspect the
graph and flow data, rerun and preferably independently reimplement the
enumerations, conduct a specialist prior-art review, and take full
responsibility for all claims.

\begin{thebibliography}{9}

\bibitem{Celmins1984}
U.~A. Celmins.
\newblock \emph{On Cubic Graphs That Do Not Have an Edge-3-Colouring}.
\newblock PhD thesis, University of Waterloo, 1984.

\bibitem{HusekSamal2026}
R.~Hu\v{s}ek and R.~\v{S}\'amal.
\newblock Exponentially many circuit double covers.
\newblock arXiv:2607.24724 [math.CO], version 1, July 27, 2026.
\newblock \url{https://arxiv.org/abs/2607.24724}.

\bibitem{Jaeger1979}
F.~Jaeger.
\newblock Flows and generalized coloring theorems in graphs.
\newblock \emph{Journal of Combinatorial Theory, Series B},
  26(2):205--216, 1979.
\newblock \url{https://doi.org/10.1016/0095-8956(79)90057-1}.

\bibitem{KralEtAl2009}
D.~Kr\'al', E.~M\'a\v{c}ajov\'a, O.~Pangr\'ac, A.~Raspaud,
  J.-S.~Sereni, and M.~\v{S}koviera.
\newblock Projective, affine, and abelian colorings of cubic graphs.
\newblock \emph{European Journal of Combinatorics}, 30:53--69, 2009.
\newblock \url{https://doi.org/10.1016/j.ejc.2007.11.029}.

\bibitem{Oum2026}
S.-i. Oum.
\newblock A proof of the Cycle Double Cover Conjecture by OpenAI:
  an exposition.
\newblock arXiv:2607.16356 [math.CO], version 2, July 23, 2026.
\newblock \url{https://arxiv.org/abs/2607.16356}.

\bibitem{Preissmann1981}
M.~Preissmann.
\newblock \emph{Sur les colorations des ar\^etes des graphes cubiques}.
\newblock PhD thesis, Universit\'e de Grenoble, 1981.

\end{thebibliography}

\end{document}
